% ============================================================
% Hesam Salmabadi — Curriculum Vitae
% Compiled with: xelatex cv_hesam_salmabadi.tex
% ATS-friendly: single column, no tables in header, clean headings
%
% CAREER TRACK INSTRUCTIONS:
%   Track A (Academic/Postdoc): keep all sections as-is; lead with Summary A
%   Track B (Geospatial Data Scientist): use Summary B; move Skills up
%   Track C (ML/Data Engineer): use Summary C; emphasise pipeline work
%   — Swap the \SUMMARY block and reorder \section{} as needed
% ============================================================
%!TEX TS-program = xelatex
\documentclass[10pt, a4paper]{article}

% --- Packages ---
\usepackage{geometry}
\geometry{top=1.8cm, bottom=1.8cm, left=2cm, right=2cm}

\usepackage[T1]{fontenc}
\usepackage[default]{lato} % Clean, modern, widely readable sans-serif

\usepackage{titlesec}
\usepackage{enumitem}
\usepackage{hyperref}
\usepackage{parskip}
\usepackage{microtype}
\usepackage{xcolor}

% --- Colours ---
\definecolor{headingblue}{HTML}{1A3A5C}   % dark navy for section headings
\definecolor{linkblue}{HTML}{1A5276}      % slightly lighter for URLs
\definecolor{lightgray}{HTML}{666666}     % dates and secondary text

% --- Hyperref setup ---
\hypersetup{
    colorlinks=true,
    urlcolor=linkblue,
    linkcolor=linkblue,
    pdfauthor={Hesam Salmabadi},
    pdftitle={Hesam Salmabadi — CV},
}

% --- Section formatting ---
\titleformat{\section}
  {\large\bfseries\color{headingblue}}
  {}{0em}{}
  [\vspace{-4pt}\textcolor{headingblue}{\rule{\linewidth}{0.6pt}}\vspace{2pt}]

\titlespacing{\section}{0pt}{10pt}{4pt}

% --- No page numbers ---
\pagestyle{empty}

% --- List spacing ---
\setlist[itemize]{noitemsep, topsep=2pt, leftmargin=1.4em}
\setlist[itemize,1]{label=\textbullet}

% --- Custom commands ---
\newcommand{\dategray}[1]{\textcolor{lightgray}{\small #1}}
\newcommand{\jobtitle}[4]{%
  \noindent\textbf{#1} \hfill \dategray{#3} \\
  \textit{#2} \hfill \dategray{#4} \\[-4pt]
}
\newcommand{\publication}[1]{\noindent #1 \vspace{3pt}}

% ============================================================
\begin{document}
% ============================================================

% --- HEADER ---
\begin{center}
  {\LARGE\bfseries\color{headingblue} Hesam Salmabadi}\\[4pt]
  {\small
    Montreal, QC, Canada \quad|\quad
    \href{mailto:hesam.salmabadi@gmail.com}{hesam.salmabadi@gmail.com} \quad|\quad
    +1 514 430 2051 \\[2pt]
    \href{https://scholar.google.com/citations?user=QAvXX3cAAAAJ}{Google Scholar} \quad|\quad
    \href{https://www.researchgate.net/profile/Hesam-Salmabadi}{ResearchGate} \quad|\quad
    \href{https://github.com/hesam-salmabadi}{GitHub} \quad|\quad
    \href{https://www.linkedin.com/in/hesam-salmabadi/}{LinkedIn}
  }
\end{center}

\vspace{4pt}

% ============================================================
% OPTIONAL PROFILE (2–3 lines max — only add when explicitly required)
% Uncomment the block that matches the target track.
% A well-structured CV speaks for itself; this section is usually unnecessary.
% ============================================================

% --- Track A: Academic / Postdoc ---
% \section{Profile}
% Environmental scientist specialising in cryosphere remote sensing and
% atmospheric dust transport. Nine first-authored publications; FRQNT doctoral
% scholar (\$75,000 CAD). PhD at UQTR (2021–2026) on non-binary soil
% freeze–thaw monitoring from passive microwave satellites.

% --- Track B: Geospatial Data Scientist ---
% \section{Profile}
% Geospatial data scientist with nine peer-reviewed publications and extensive
% experience in Python-based pipelines for continental-scale satellite datasets
% (xarray, dask, rasterio). Currently at Agriculture and Agri-Food Canada
% building national LULC modelling frameworks; contributing author,
% Canada's Changing Climate Report 2026.

% --- Track C: ML / Data Engineer ---
% \section{Profile}
% Data scientist with a research background in large-scale geospatial ML
% pipelines. Experienced in model architecture trade-offs, spatial
% cross-validation design, and HDF5/NetCDF data engineering at scale.
% Currently at AAFC; working toward cloud deployment of geospatial models.

% ============================================================
\section{Experience}
% ============================================================

\jobtitle{Geospatial Data Scientist}
         {Agriculture and Agri-Food Canada (AAFC), Earth Observations, Science \& Technology Branch}
         {May 2025 – Present}
         {Ottawa, ON, Canada (Remote)}
\begin{itemize}
  \item Assembled a national LULC data cube at 100\,m resolution from
        multi-source satellite time series and environmental covariates
        covering all of Canada south of 60°N; engineered out-of-core
        ingestion and harmonisation pipelines (xarray, dask, Zarr)
        to support continental-scale spatial analysis
  \item Validated satellite-derived agricultural land estimates against
        Statistics Canada Census data across all 10 southern provinces;
        identified three distinct regimes of agricultural change ---
        active conversion to urban land, gradual abandonment on marginal
        land, and administrative flux --- with strong EO–Census agreement
        (Pearson $r = 0.988$)
  \item Developed Bayesian provincial demand forecasting models (PyMC)
        for LULC trajectories, reporting posterior medians and 94\% HDI
        to quantify projection uncertainty; translated posteriors into
        five policy-relevant land-use demand scenarios
  \item Built pixel-level logistic regression suitability classifiers
        with geographic block cross-validation to prevent spatial leakage
        across disjoint test regions
        (ROC-AUC 0.957, Recall 0.81 on geographically held-out blocks)
  \item Ran Dyna-CLUE with demand scenarios and pixel-level suitability
        surfaces to generate spatially explicit LULC projection maps
        across Canada
\end{itemize}

\vspace{6pt}

\jobtitle{PhD Researcher}
         {Université du Québec à Trois-Rivières (UQTR) \& Centre d'Études Nordiques, Université Laval}
         {May 2021 – May 2026}
         {Trois-Rivières, QC, Canada}
\begin{itemize}
  \item Recast the Soil Freezing Characteristic Curve (SFCC) into
        permittivity–temperature space, eliminating moisture-specific
        calibration; applied Bayesian hierarchical partial pooling across
        87 in situ stations (10 networks, boreal, prairie, and tundra) —
        mean R$^{2}$ = 0.95; demonstrated that the zero-curtain transitional
        state persists $>$100 days in moisture-rich boreal sites, invisible
        to standard binary classifiers
  \item Developed a probabilistic non-binary freeze–thaw retrieval fusing
        SMAP (L-band) and AMSR2 passive microwave observations into a
        continuous Frozenness Index and 3-class state map at 9 km resolution
        across the Northern Hemisphere (10-year record); built end-to-end
        ingestion and continental-scale processing pipelines via Dask and
        Zarr; validated against a three-product unanimous consensus ---
        match rates 78.8--98.2\%, Cohen's Kappa 0.44--0.87, PR AUC $>$0.95
        in three open biomes; applied DBSCAN clustering and
        StratifiedGroupKFold to eliminate spatial autocorrelation leakage
  \item Applied XGBoost with SHAP attribution over a 9-year continental
        record; identified sub-canopy snowpack as the primary driver of
        frozen ground duration in forested biomes --- a regime shift
        invisible to temperature-only models (R$^{2}$ = 0.894,
        MAE = 11.5 days); linked the Frozenness Index to AmeriFlux
        ecosystem respiration via piecewise regression, confirming the
        zero-curtain onset marks the metabolic shutdown tipping point for
        cold-season carbon flux
  \item Conducted field campaigns across remote Canadian ecosystems;
        responsible for sensor installation, data logger servicing, and
        in-field quality control across boreal, subarctic, and prairie terrain
  \item Funded by FRQNT Doctoral Research Scholarship (\$75{,}000 CAD,
        2023--2026), UQTR Bourse Universalis Causa (\$2{,}100 CAD), and
        Canadian Space Agency Research Grant ($\sim$\$30{,}000 CAD via
        supervisor)
\end{itemize}

\vspace{6pt}

\jobtitle{Geospatial Data Analyst — Contributing Author}
         {Environment and Climate Change Canada (ECCC)}
         {Apr 2024 – Oct 2024}
         {Toronto, ON, Canada (Remote)}
\begin{itemize}
  \item Developed a multi-decade satellite data pipeline processing 40+
        years of daily freeze–thaw observations into structured xarray
        spatial time series across Canada, underpinning the national-scale
        cryosphere trend assessment
  \item Quantified a reduction of 1.7 days/decade in frozen ground days
        since 1979 and a 28\% increase in unfrozen days across the Canadian
        Tundra (unfrozen season lengthened by $\sim$27 days); produced
        publication-ready spatial trend maps for the federal climate assessment
  \item Contributing author, \textit{Canada's Changing Climate Report 2026},
        Chapter 6: Cryosphere Changes (forthcoming)
\end{itemize}

\vspace{6pt}

\jobtitle{Field Research Technician}
         {NASA Jet Propulsion Laboratory (JPL) — SMAPVEX22-Boreal Campaign}
         {Aug 2022}
         {Saskatchewan, Canada (On-site)}
\begin{itemize}
  \item Contributed to ground-based soil moisture and temperature measurements
        across the BERMS study area as part of a NASA SMAP satellite validation
        campaign; field data are part of Berg et al.\ (2025), \textit{IEEE JSTARS}
\end{itemize}

\vspace{6pt}

\jobtitle{Environmental Remote Sensing Scientist}
         {Research Institute of Cultural Heritage \& Tourism (RICHT)}
         {Oct 2019 – Jan 2021}
         {Tehran, Iran}
\begin{itemize}
  \item Designed an automated Lagrangian simulation pipeline (Python/PySPLIT)
        executing forward trajectories over 18 years to map
        altitude-dependent transport of salt and dust sediments from
        Lake Urmia's desiccating playa
  \item Engineered a multi-decadal satellite processing pipeline (MODIS AOD,
        Landsat) to track lake desiccation and identify active saline dust
        sources; established a strong correlation ($R = 0.86$) between
        exposed lake bed area and regional aerosol events
  \item Quantified a 76\% loss in lake surface area linked to agricultural
        expansion; demonstrated that low-altitude dust transport drove a
        30--60\% increase in dust hours for downwind cities
\end{itemize}

\vspace{6pt}

\jobtitle{MSc Researcher}
         {Iran University of Science and Technology (IUST)}
         {Sep 2015 – Jul 2018}
         {Tehran, Iran}
\begin{itemize}
  \item Investigated the air quality impact of wetland desiccation in the
        Middle East; combined multi-decadal Landsat time-series for wetland
        area monitoring with HYSPLIT forward trajectory modelling --- identified
        four primary dust transport pathways from the Mesopotamian Marshes,
        with regional implications for air quality from the Persian Gulf to
        Central Asia
  \item Monitored national SO$_2$ column concentrations over Iran
        (2005--2016) using Aura/OMI satellite retrievals; characterised
        seasonal transport corridors from an industrial point source via
        HYSPLIT forward trajectories
  \item Contributed geospatial analysis to a collaborative environmental
        study on heavy metal (Hg) distribution in Musa Estuary; produced
        spatial contamination maps using ArcGIS and geostatistical
        interpolation
  \item Produced 5 peer-reviewed publications (2019--2025); graduated as
        top-ranked MSc student (GPA 18.33/20, national scholarship)
\end{itemize}

% ============================================================
\section{Education}
% ============================================================

\jobtitle{Doctor of Philosophy — Environmental Sciences}
         {Université du Québec à Trois-Rivières (UQTR)}
         {May 2021 – May 2026}
         {Trois-Rivières, QC, Canada}

\noindent\textit{Thesis:} Non-Binary Monitoring of Seasonally Frozen Ground:
From In Situ Sensors to Continental Passive Microwave Retrievals\\
\textit{Supervisor:} Prof.\ Alexandre Roy

\vspace{6pt}

\jobtitle{MSc — Civil Engineering (Environmental Engineering)}
         {Iran University of Science and Technology (IUST)}
         {Sep 2015 – Jul 2018}
         {Tehran, Iran}

\noindent\textit{Thesis:} Determining the Effect of Wetland and Lake
Desiccation in Iran on the Air Quality of Potentially Affected Regions\\
\textit{GPA:} 18.33/20 — First Ranked Student \quad
\textit{Supervisor:} Prof.\ Mohsen Saeedi

\vspace{6pt}

\jobtitle{BSc — Civil Engineering}
         {University of Birjand}
         {Sep 2010 – Sep 2014}
         {Birjand, Iran}

% ============================================================
\section{Selected Publications}
% ============================================================

\noindent\textcolor{lightgray}{\small
  14 publications total (9 first-authored) \,|\,
  Citations: 181 \,|\, h-index: 6 \,|\, i10-index: 5 \quad
  \href{https://scholar.google.com/citations?user=QAvXX3cAAAAJ}{[Google Scholar]}
}

\vspace{4pt}

% --- Swap / reorder entries to match the target role ---

\publication{\textbf{Salmabadi H.}, Pardo Lara R., Berg A., Mavrovic A., Hanes C., Montpetit B., Roy A. (2026).
In situ monitoring of seasonally frozen ground using soil freezing characteristic curve in permittivity–temperature space.
\textit{The Cryosphere}, 20, 1635–1654.
\href{https://doi.org/10.5194/tc-20-1635-2026}{doi:10.5194/tc-20-1635-2026}}

\publication{\textbf{Salmabadi H.}, Berg A., Roy A. (2026, preprint).
Non-Binary Retrievals of Soil Freeze–Thaw Dynamics Using Multi-Frequency Passive Microwave Data.
\textit{SSRN}. Under review.
\href{https://ssrn.com/abstract=6626104}{ssrn.com/abstract=6626104}}

\publication{\textbf{Salmabadi H.}, Saeedi M., Notaro M., Roy A. (2025).
Dust transport pathways from the Mesopotamian Marshes.
\textit{Aeolian Research}, 73, 100975.
\href{https://doi.org/10.1016/j.aeolia.2025.100975}{doi:10.1016/j.aeolia.2025.100975}}

\publication{\textbf{Salmabadi H.}, Saeedi M., Roy A., Kaskaoutis D.G. (2023).
Quantifying the contribution of Middle Eastern dust sources to PM\textsubscript{10}
levels in Ahvaz, Southwest Iran.
\textit{Atmospheric Research}, 295, 106993.
\href{https://doi.org/10.1016/j.atmosres.2023.106993}{doi:10.1016/j.atmosres.2023.106993}}

\publication{\textbf{Salmabadi H.}, Roy A., Berg A. (2026).
Capturing the Soil Zero-Curtain Effect from Multi-Frequency Passive Microwave Retrievals.
\textit{ISPRS Congress 2026}.}

\publication{Berg A.A., \ldots, \textbf{Salmabadi H.}, \ldots (2025).
Soil Moisture Active/Passive Validation Experiment Within the Canadian Boreal Forest in 2022 (SMAPVEX22-Boreal).
\textit{IEEE J. Sel. Topics Appl. Earth Obs. Remote Sens.}, 18, 11816–.
\href{https://doi.org/10.1109/JSTARS.2025.3564195}{doi:10.1109/JSTARS.2025.3564195}}

% ============================================================
\section{Skills}
% ============================================================

\noindent\textbf{Programming \& Data:} Python (xarray, dask, numpy, pandas,
rasterio, geopandas, scikit-learn, matplotlib, cartopy, h5py), GDAL (Python
\& CLI), Bash, MATLAB, LaTeX, Git/GitHub, MySQL

\vspace{3pt}
\noindent\textbf{Machine Learning \& Statistics:} Logistic regression (L1/L2),
random forests, SVM, DBSCAN, Bayesian hierarchical partial pooling,
probabilistic classification, Mann-Kendall trend test, spatial cross-validation

\vspace{3pt}
\noindent\textbf{Remote Sensing:} SMAP (L-band), AMSR2, MODIS AOD,
Landsat time-series, OMI SO\textsubscript{2}, ERA5, GDAS;
passive microwave freeze–thaw retrieval; source apportionment (PSCF)

\vspace{3pt}
\noindent\textbf{Modelling:} HYSPLIT trajectory modelling (forward \&
backward); numerical atmospheric transport; soil freeze–thaw modelling

\vspace{3pt}
\noindent\textbf{Data Formats:} NetCDF, HDF4/5, GeoTIFF, Zarr, Shapefile

\vspace{3pt}
\noindent\textbf{GIS:} QGIS, ArcGIS/ArcMap, PostGIS (learning), spatial
interpolation (IDW, kriging concepts)

\vspace{3pt}
\noindent\textbf{Cloud \& Infra:} AWS (in progress), GitHub Actions basics,
agentic AI tooling

% ============================================================
\section{Awards \& Funding}
% ============================================================

\noindent\textbf{FRQNT Doctoral Research Scholarship} \hfill \dategray{2023–2026}\\
Fonds de recherche du Québec – Nature et technologies \quad \$75,000 CAD\\
Project: \textit{Development of a New Passive Microwave Multi-Sensor Freeze/Thaw
Data Product to Improve Growing Season Length Monitoring over North America}

\vspace{4pt}
\noindent\textbf{Bourse Universalis Causa} \hfill \dategray{2021–2024}\\
Université du Québec à Trois-Rivières \quad Up to \$21,000 CAD + tuition exemption

\vspace{4pt}
\noindent\textbf{First Ranked Student — MSc Program} \hfill \dategray{2018}\\
Iran University of Science and Technology (IUST) \quad GPA: 18.33/20

\vspace{4pt}
\noindent\textbf{Full Tuition Scholarship — BSc \& MSc} \hfill \dategray{2010–2018}\\
Ministry of Science, Research and Technology of Iran — awarded on national
competitive examination ranking

% ============================================================
\section{Certificates \& Professional Development}
% ============================================================

\begin{itemize}
  \item AI for Earth Program — Amii (Alberta Machine Intelligence Institute), 2026
  \item Advanced Learning Algorithms — DeepLearning.AI / Coursera, 2025
  \item Supervised Machine Learning: Regression and Classification — DeepLearning.AI / Coursera, 2025
  \item Google Data Analytics Certificate (Foundations, Python, Data Exploration) — Google / Coursera, 2022–2025
  \item Advanced Python Programming — Maktabkhooneh, 2020
\end{itemize}

% ============================================================
\section{Languages}
% ============================================================

Persian (Native) \quad|\quad English (Fluent)

% ============================================================
\end{document}
